\chapter{Extension fields: $k=4$ and the $D_4$ symmetry}\label{ch:extension}

Over $\FF_{q^k}$ with a factor base of points whose abscissa lies in
$\FF_q$, Gaudry's index calculus~\cite{gaudry2009} decomposes each target
into $k$ summands by Weil restriction of $\Sem{k+1}$. It costs
$\tilde O(q^{2-2/k})$, below rho's $q^{k/2}$ for $k\ge3$, but with a
Gröbner constant that grows very fast in $k$. Faugère, Gaudry, Huot and
Renault~\cite{fghr2014} show that on curves with rational 2-torsion the
system has the larger symmetry group $D_k$, the Weyl group of type D, in
a well-chosen coordinate. At $k=4$ this chapter reports how far that
symmetry moves the measured gap to rho. Chapter~\ref{sec:prime}
(\S\ref{sec:gaudry}) gives the odd-characteristic engineering in detail.

\section{The $D_4$ system}

On $E$ with rational 2-torsion point $T$, the coordinate
$Y=(x-x_0-c)/(x-x_0+c)$ satisfies $Y(P+T)=-Y(P)$. In $Y$, $\Sem5$ is
invariant under $D_4$, of order 192. Written in the invariants
$\sigma_k=e_k(Y_i^2)$ and $\pi=\prod Y_i$, the system shrinks:
\begin{itemize}
\item the Bézout number falls from 4{,}096 to 512;
\item the Macaulay degree falls from 29 to 14;
\item the matrix falls from $41{,}780\times40{,}920$ to $13{,}412\times11{,}628$.
\end{itemize}
The engineering ladder, each step against the previous incumbent, was:
\begin{enumerate}
\item \emph{$D_4$ invariants.} One decomposition costs $C_4=2.54\cdot10^{10}$
  $\FF_p$ multiplications instead of $1.21\cdot10^{12}$ at $p=269$, a
  factor of $47.5$--$48.4$. All 19 solved targets agree with an independent
  meet-in-the-middle enumeration
  (\art{research/gaudry_quartic_d4_20261013/}).
\item \emph{The $Y$-line on a complete ladder.} On $p=101$--$1033$ (21
  curves) the cost is $0.0548$ $[0.048,0.061]$ of plain $\Sem4$, at the same
  exponent.
\item \emph{Sparsest pivots, the weighted grading $\sigma_4=\pi^2$, and shape
  replay.} $C_4$ falls to $6.89\cdot10^7$, by $369.8\times$
  $[367.9,371.6]$; all 648 of 648 solves are exact. Ecbench $S$ falls
  $398$--$406\times$ (\art{research/gaudry_quartic_d4_cost_20261014/}).
\item \emph{First-dependency stop.} The cost is $0.979$ $[0.973,0.986]$ of the
  previous step. Double large primes cut the linear algebra to $0.827\times$
  but regress the whole solve, to $1.025\times$.
\item \emph{A Frobenius-reduced form in characteristic two.} Weights
  $(2,2,1,1)$ give Macaulay degree 11 and a $384\times448$ matrix. One
  decomposition is $481\times$ $[479,483]$ below the plain form at $m=8$
  (\art{research/gaudry_quartic_d4_char2_20261015/}).
\end{enumerate}
The $D_n$ symmetry is FGHR's~\cite{fghr2014}. Invariant rings in
characteristic two, including even-power substitutions, were studied by
Galbraith and Gebregiyorgis~\cite{galbraith2014}. Whether the
Frobenius-reduced weighting is in that work has not been checked.

\section{Distance to rho}

\paragraph{Odd characteristic.} Over $\FF_{p^4}$, index calculus is still
$1.39\cdot10^6\times$ rho at the measured sizes. The modelled handover
moves from $2^{150}$ to about $2^{103}$--$2^{119}$, depending on the frame,
but that is an extrapolation through a measured constant.

\paragraph{Characteristic two.} Over $\FF_{(2^m)^4}$ the gap closes with
size. IC over the negation-map rho model is:

\smallskip
\begin{tabular}{@{}lrrrrr@{}}
\toprule
$m$ & 6 & 8 & 10 & 12 & 14 \\
$\log_2 n$ & 24 & 32 & 40 & 48 & 56 \\
\midrule
IC / rho model & 5{,}080 & 648 & 80.6 & 49.7 & 6.06 \\
IC / measured rho & 4{,}427 & 919 & 98.9 & 47.2 & 4.27 \\
\bottomrule
\end{tabular}
\smallskip

\noindent These are first targets; $m=14$ is a single target. The fitted
slope is $-0.248$ in $\log_2 n$, which is Gaudry's $n^{-1/4}$. The
extrapolated crossover was $n\approx2^{68}$, where rho takes about $2^{34}$
steps, so it could be tested directly.

\paragraph{The red-team verdict.} An adversarial review tested that
crossover (\art{research/gaudry_k4_char2_redteam_20261026/REPORT.md}). Its
verdict: \textbf{not survived (unsupported), not refuted}.
\begin{itemize}
\item \emph{Data.} The only clean point is at $m=15$, with two targets per
  arm. The recorded IC/rho there is $6.30$. After the corrections that
  favour index calculus it is $5.2$, with range $2.4$--$15.0$ over
  unit-price variants. The $m=17$ ratio of $0.72$ came from a killed run
  and is not evidence. No interval is possible with $n=2$.
\item \emph{Rho reference survives.} These are generic ordinary curves with
  $\operatorname{Aut}=\{\pm1\}$, and the floor $0.886$ is about 3.5\%
  generous to rho.
\item \emph{Pricing partly survives.} The unit price moves the ratio by
  $0.64$--$0.89$, with a plausible range of $0.31$--$3.1$. Unpriced rho
  work moves it by about $0.93$.
\item \emph{Fitted exponents} (cost in $q=2^m$): IC relations $1.07$
  $[1.01,1.13]$, IC linear algebra $2.03$, rho $2.00$ $[1.71,2.29]$. A
  prime-field cross-check agrees.
\item \emph{Corrected crossover.} The IC/rho ratio falls by about
  $2^{-0.93}$ per unit of $m$, which puts the crossover at
  $m^*\approx17.6$, with range $16.3$--$19.2$: at or above the edge of what
  was measured.
\end{itemize}
Linear algebra is 0.4\% of the cost at $m=15$. Basic Gaudry linear
algebra grows as $q^2$, like rho, so the $n^{-1/4}$ trend lasts only while
the relation phase dominates. Beyond that point double large primes are
required.

The red team froze a plan to settle the question: $m=14$--$17$, eight fresh
targets per arm, a strong rho with batched inversion and charged
canonicalisation, counted and hardware units, about 30\,h of wall time.
It has not run. This is the cheapest experiment in the programme that
could yield a measured crossover of index calculus below rho on an
elliptic curve. It would be a crossover in a weak family:
$\FF_{(2^m)^4}$ is not a field deployed curves use.

\section{Other extension-field routes}

\begin{itemize}
\item \emph{Joux--Vitse $(n-1)$-decomposition.} The
  $(n-1)$-decomposition~\cite{jouxvitse2012} at $n=5$ over $E(\FF_{p^5})$
  completed verified logarithms at $4.5\cdot10^7$--$1.8\cdot10^8\times$ rho,
  with a crossover extrapolated to groups near $2^{291}$
  (\art{research/joux_vitse_n1_20261015/}).
\item \emph{Joux--Vitse cover.} The cover of $E(\FF_{p^6})$ is in
  Chapter~\ref{sec:transfer}.
\end{itemize}
